\documentclass[10pt,a4paper]{article}
\usepackage[margin=18mm,top=20mm,bottom=19mm,headheight=14pt]{geometry}
\usepackage[T1]{fontenc}
\usepackage{lmodern}
\usepackage{microtype}
\usepackage{xcolor}
\usepackage{tikz}
\usetikzlibrary{arrows.meta,positioning,shapes.geometric}
\usepackage{listings}
\usepackage{tabularx,booktabs,array}
\usepackage{enumitem}
\usepackage{fancyhdr}
\usepackage{titlesec}
\usepackage{hyperref}
\definecolor{ink}{HTML}{142D40}
\definecolor{teal}{HTML}{007F82}
\definecolor{muted}{HTML}{536575}
\definecolor{pale}{HTML}{F0F6F8}
\definecolor{rulegray}{HTML}{D6E1E7}
\definecolor{red}{HTML}{B9424A}
\hypersetup{colorlinks=true,linkcolor=teal,urlcolor=teal}
\renewcommand{\familydefault}{\sfdefault}
\color{ink}
\setlength{\parindent}{0pt}
\setlength{\parskip}{5pt}
\setlength{\emergencystretch}{2em}
\renewcommand{\arraystretch}{1.2}
\setlist[itemize]{leftmargin=16pt,itemsep=2pt,topsep=2pt}
\setlist[enumerate]{leftmargin=18pt,itemsep=2pt,topsep=2pt}
\titleformat{\section}{\Large\bfseries\color{ink}}{\textcolor{teal}{\thesection\enspace}}{0pt}{}
\titleformat{\subsection}{\normalsize\bfseries\color{teal}}{}{0pt}{}
\titlespacing*{\section}{0pt}{9pt}{7pt}
\titlespacing*{\subsection}{0pt}{7pt}{3pt}
\pagestyle{fancy}
\fancyhf{}
\fancyhead[L]{\small\textcolor{muted}{NESTJS TP1}}
\fancyhead[R]{\small\textcolor{muted}{Study documentation}}
\fancyfoot[L]{\footnotesize\textcolor{muted}{Explain the responsibility, then the code.}}
\fancyfoot[R]{\small\thepage}
\renewcommand{\headrulewidth}{0.4pt}
\lstdefinelanguage{TypeScript}{keywords={import,from,export,class,const,let,return,new,private,readonly,constructor,async,await,if,throw,extends,string,number,boolean,void,true,false,enum},sensitive=true,comment=[l]{//},morecomment=[s]{/*}{*/},morestring=[b]',morestring=[b]"}
\lstset{language=TypeScript,basicstyle=\ttfamily\footnotesize,keywordstyle=\color{teal}\bfseries,commentstyle=\color{muted},stringstyle=\color{red},backgroundcolor=\color{pale},frame=single,rulecolor=\color{rulegray},framerule=0.4pt,framesep=6pt,xleftmargin=6pt,xrightmargin=6pt,breaklines=true,columns=fullflexible,keepspaces=true,showstringspaces=false,tabsize=2,aboveskip=7pt,belowskip=6pt}
\newcommand{\code}[1]{\texttt{#1}}
\newcommand{\note}[1]{\par\smallskip\noindent\colorbox{pale}{\parbox{\dimexpr\linewidth-12pt}{\vspace{4pt}\textcolor{teal}{\textbf{Remember}} \ #1\vspace{4pt}}}\par\smallskip}
\tikzset{box/.style={draw=rulegray,fill=pale,rounded corners=3pt,align=center,minimum height=11mm,text width=38mm,font=\small},accent/.style={box,draw=teal,fill=teal!7},arr/.style={-{Latex[length=2mm]},thick,draw=teal}}
\begin{document}


\hypersetup{pdftitle={TP1 NestJS Handoff},pdfauthor={Yassine}}
\thispagestyle{empty}
{\small\bfseries\color{teal} NESTJS / RT4 TP1 / HANDOFF FOR THE NEXT DEVELOPER}\par\vspace{10mm}
{\fontsize{31}{35}\selectfont\bfseries TP1 Continuation Notes\par}\vspace{4mm}
{\Large Pull the repository, verify the baseline, continue from slide 6\par}
\vspace{5mm}{\color{muted} Project: \code{todo-api} \textbullet\ Current boundary: TypeORM configured}\par\vspace{8mm}
This handoff records exactly what is implemented and gives a safe continuation path. Do not start by rewriting the existing architecture. Run the baseline checks first, then extend it feature by feature.
\note{The repository is intentionally stopped before validation and Todo CRUD. Your first task is slide 6: DTO constraints and validation.}

\section{Clone and verify}
After the repository is published, clone it with:
\begin{lstlisting}[language={}]
git clone https://github.com/MohamedYassineTurki/tp1-nestjs-todo-api.git
cd tp1-nestjs-todo-api
npm install
npm run build
\end{lstlisting}
Start development mode:
\begin{lstlisting}[language={}]
npm run start:dev
\end{lstlisting}
In another terminal, verify the completed endpoint:
\begin{lstlisting}[language={}]
curl http://localhost:3000/uuid-test
\end{lstlisting}
Expected response:
\begin{lstlisting}[language={}]
{"id":"a-generated-uuid"}
\end{lstlisting}
Repeat the request. The value should change.

\section{Current structure}
\begin{lstlisting}[language={}]
src/
  main.ts                         Application startup
  app.module.ts                   Root module and TypeORM connection
  app.controller.ts               Initial Hello World route
  app.service.ts                  Initial service
  common/
    common.module.ts              Global shared module
    constants/common.constants.ts UUID provider token
  uuid-test/
    uuid-test.module.ts           Test feature module
    uuid-test.controller.ts       GET /uuid-test
    uuid-test.service.ts          Injects and calls UUID provider
  todos/
    entities/todo.entity.ts       TypeORM Todo entity
    enums/status.enum.ts          Temporary three-value status enum
todo.sqlite                       Local development database
\end{lstlisting}
The project currently has no Todo controller or Todo service. The entity exists, but no Todo repository is injected yet.

\section{Completed work}
\subsection{Nest application}
\code{src/main.ts} bootstraps \code{AppModule} and listens on port 3000. \code{AppModule} registers the root controller, root service, common module, UUID test module, and TypeORM connection.

\subsection{UUID provider}
\code{CommonModule} imports \code{v4 as uuidv4} from \code{uuid}, registers it with \code{useValue}, and exports the symbol token \code{UUID\_PROVIDER}. The global decorator makes the shared provider available to feature modules.
\begin{lstlisting}
{
  provide: UUID_PROVIDER,
  useValue: uuidv4,
}
\end{lstlisting}
\code{UuidTestService} receives the function with \code{@Inject(UUID\_PROVIDER)}. Its controller exposes \code{GET /uuid-test}.

\subsection{Entity and database}
\code{TodoEntity} contains an auto-generated integer ID, name, description, creation timestamp, and status. \code{AppModule} configures TypeORM with \code{better-sqlite3}, the file \code{todo.sqlite}, and \code{synchronize: true} for local learning.
The earlier configuration \code{type: 'sqlite'} failed with the installed TypeORM definitions. The working driver name is \code{better-sqlite3}; both the package and the configuration are already present.

\section{Your first continuation: slide 6}
Install validation packages:
\begin{lstlisting}[language={}]
npm install class-validator class-transformer
\end{lstlisting}
Add the global validation pipe in \code{src/main.ts}:
\begin{lstlisting}
import { ValidationPipe } from '@nestjs/common';

app.useGlobalPipes(
  new ValidationPipe({
    whitelist: true,
    forbidNonWhitelisted: true,
  }),
);
\end{lstlisting}
Create \code{src/todos/dto/create-todo.dto.ts}. The assignment requires a required description of at least 10 characters and a required name between 3 and 10 characters:
\begin{lstlisting}
import { IsString, Length, MinLength } from 'class-validator';

export class CreateTodoDto {
  @IsString()
  @Length(3, 10)
  name!: string;

  @IsString()
  @MinLength(10)
  description!: string;
}
\end{lstlisting}
Create a separate update DTO with optional fields, retaining the size constraints. Add status validation with \code{IsEnum(StatusEnum)}. Put all custom error messages in one reusable constants file.

\section{Remaining TP sequence}
Follow this order so each step has a working foundation:
\begin{enumerate}
\item Slide 6: validation pipe, create DTO, update DTO, reusable messages, enum validation.
\item Slide 7: add \code{updatedAt} and \code{deletedAt}; keep \code{createdAt} controlled by TypeORM and excluded from update input.
\item Slide 8: extract creation, update, and deletion timestamp columns into a reusable base entity.
\item Slide 9: generate \code{TodosModule}, \code{TodosService}, and \code{TodosController};\\ import \code{TypeOrmModule.forFeature([TodoEntity])}; implement \code{addTodo} and preserve both requested versions.
\item Slide 10: implement update by ID and handle a missing Todo.
\item Slides 11-12: implement hard delete, soft delete, and restore.
\item Slide 13: add one aggregation endpoint returning counts for the three statuses.
\item Slides 14-15: get all Todos and get one Todo by ID with \code{NotFoundException}.
\item Slide 16: add optional search and status filters with the intended \code{(name OR description) AND status} logic.
\item Slide 17: add validated page and limit parameters, stable ordering, skip, take, and useful pagination metadata.
\end{enumerate}

\section{Recommended Todo wiring}
When you create the feature module, the repository registration belongs in that module:
\begin{lstlisting}
@Module({
  imports: [TypeOrmModule.forFeature([TodoEntity])],
  controllers: [TodosController],
  providers: [TodosService],
})
export class TodosModule {}
\end{lstlisting}
The service then receives the repository:
\begin{lstlisting}
constructor(
  @InjectRepository(TodoEntity)
  private readonly todoRepository: Repository<TodoEntity>,
) {}
\end{lstlisting}
Keep HTTP extraction in the controller and database/application logic in the service.

\section{Baseline issues to know}
The application build passes and the UUID endpoint works. The CLI-generated unit tests for \code{UuidTestService} and \code{UuidTestController} currently fail because their isolated test modules do not provide the UUID token and service dependency. Before final delivery, update those tests with mocks or imported modules. Do not mistake a unit-test dependency error for a runtime application error.
The local file \code{todo.sqlite} should not be committed. Add it to \code{.gitignore}; the database can be recreated from the entity and TypeORM configuration. Also ignore \code{node\_modules}, \code{dist}, and \code{coverage}.

\section{Git workflow}
After each coherent TP milestone:
\begin{lstlisting}[language={}]
git status
git add src package.json package-lock.json .gitignore
git commit -m "Implement slide 6 validation"
git push
\end{lstlisting}
Do not commit passwords, database credentials, or generated dependency folders. Keep commits small enough that the professor can follow the architecture.

\section{Slides 10-12: update, delete, soft delete, restore}
This part was added by Zak on the branch \code{zak}, created from \code{eya}. Its responsibility is to change or remove an existing Todo by its ID. Everything is in \code{src/todos/todos.controller.ts}. It follows the same approach as \code{addTodo}: the controller uses the injected \code{DataSource} and calls \code{getRepository(TodoEntity)}. There is still no \code{TodosService} or \code{TodosModule} (slide 9), so when slide 9 is done the bodies of these four methods can move into the service.

\subsection{Routes}
\begin{itemize}
\item Slide 10, \code{PATCH /todos/:id}, method \code{updateTodo}: calls \code{repository.update(Number(id), updateTodoDto)}. Only the fields sent in the body change, and the DTO validation still applies (a name such as \code{ab} returns 400).
\item Slide 11, \code{DELETE /todos/:id}, method \code{deleteTodo}: calls \code{repository.delete}, which removes the row for good.
\item Slide 12, \code{DELETE /todos/:id/soft}, method \code{softDeleteTodo}: calls \code{repository.softDelete}, which only fills \code{deletedAt} (the \code{@DeleteDateColumn} of \code{TimestampEntity}). The row stays in the table.
\item Slide 12, \code{PATCH /todos/:id/restore}, method \code{restoreTodo}: calls \code{repository.restore}, which sets \code{deletedAt} back to null.
\end{itemize}
The ID arrives from the URL as text, so \code{Number(id)} converts it before it is passed to TypeORM.


\subsection{Verification}
Tested in PowerShell with \code{Invoke-RestMethod} while \code{npm run start:dev} was running: creating a Todo with \code{POST /v2/todos}, updating its status, updating its name, soft delete, restore and hard delete all returned their success message. Updating with the name \code{ab} was rejected with a 400 and the custom validation message.
\begin{lstlisting}[language={}]
Invoke-RestMethod -Method Patch -Uri http://localhost:3000/todos/1 -ContentType "application/json" -Body '{"status":"IN_PROGRESS"}'
Invoke-RestMethod -Method Delete -Uri http://localhost:3000/todos/1/soft
Invoke-RestMethod -Method Patch -Uri http://localhost:3000/todos/1/restore
Invoke-RestMethod -Method Delete -Uri http://localhost:3000/todos/1
\end{lstlisting}
\note{Known limitation: this is the minimal version with no existence check. If the ID does not exist, the methods change nothing but still return the success message. Slide 10 asks to handle a missing Todo, so this is still to do: for example call \code{findOne} first and throw \code{NotFoundException}.}


\end{document}